\chapter{Cerebrospinal-Fluid Oligoclonal Bands}

\section*{What Is This Test?}
CSF oligoclonal-band testing identifies restricted immunoglobulin bands produced within the central nervous system. The bands are detected by comparing cerebrospinal fluid with a paired serum sample using isoelectric focusing and immunoblotting. They indicate intrathecal immunoglobulin synthesis but are not specific to one disease.

\section*{Alternative Names}
Oligoclonal Bands; CSF OCB; OCB Testing; Intrathecal IgG Bands; CSF Immunofixation

\section*{Why It Is Ordered}
\begin{itemize}
  \item Supporting evaluation of suspected multiple sclerosis or another inflammatory demyelinating disease
  \item Investigating chronic inflammatory or infectious disorders of the central nervous system
  \item Complementing CSF cell count, protein, IgG index, and other neurologic testing
\end{itemize}

\section*{How It Is Performed}
Cerebrospinal fluid and serum collected as close together as possible are tested in parallel. Isoelectric focusing separates immunoglobulins by charge, and immunoblotting detects the resulting IgG bands. A paired serum sample is essential for distinguishing bands produced within the central nervous system from bands that are present systemically.

\section*{Interpretation and Limitations}
Two or more CSF-restricted bands, absent from the paired serum, support intrathecal IgG synthesis. Positive bands can occur in multiple sclerosis, infections, autoimmune disease, and some malignancies. A negative result does not exclude multiple sclerosis or other inflammatory disease. Interpretation requires the clinical presentation, neurologic examination, MRI findings, and the remainder of the CSF profile; lumbar puncture is a specimen-collection procedure outside this chapter.
\section*{Regulatory Status and Authorization Date}
Regulatory status applies to the specific assay, instrument, manufacturer, and intended use rather than to the test name alone. No single approval date can be assigned to this test category. For a commercial assay, verify the current product in the relevant regulator's database: the U.S. Food and Drug Administration (FDA) clearance, approval, or authorization; India's Central Drugs Standard Control Organisation (CDSCO) licence or permission; the European Union In Vitro Diagnostic Regulation (EU IVDR) CE certificate issued through the applicable conformity-assessment route; or the United Kingdom Medicines and Healthcare products Regulatory Agency (MHRA) registration and UKCA/CE status. Laboratory-developed tests may not have an individual FDA, CDSCO, EU IVDR, or MHRA approval date.
\clearpage
